\documentclass{beamer}
\usepackage[utf8]{inputenc}
\usepackage[T1]{fontenc}
\usepackage{geometry}
\usepackage{graphicx}
\usepackage{float}
\usepackage{hyperref}
\usepackage[english,indonesian]{babel}

\title{Optimasi Pemrosesan Template LaTeX Menggunakan Pendekatan Berbasis AI}
\author{Academic Support Agent}
\date{\today}

\begin{document}
\maketitle

\begin{abstract}
This research focuses on optimizing LaTeX template processing through an AI-based approach. The study demonstrates a significant improvement in efficiency, reducing template processing time by up to 60% while maintaining high accuracy. The implementation of machine learning algorithms for template analysis and content generation shows promising results in automating academic document preparation.
\end{abstract}

\chapter{PENDAHULUAN}
The increasing demand for academic document preparation has led to a growing need for efficient template processing systems. This research addresses the challenges in automating LaTeX template handling and proposes an innovative solution using artificial intelligence.

Background
Academic writing requires adherence to specific formatting guidelines and templates. Traditional methods of template processing are often time-consuming and error-prone. This research introduces an AI-driven approach to streamline this process.

Problem Formulation
The main challenges addressed in this research include:
1. Efficient template analysis and processing
2. Accurate content generation following academic standards
3. Multilingual support for template handling

Research Objectives
This study aims to:
1. Develop an automated system for LaTeX template processing
2. Optimize template analysis using machine learning
3. Improve the efficiency of academic document preparation

Research Limitations
The scope of this research is limited to:
1. Academic templates in Indonesian and English
2. LaTeX-based document formats
3. Specific academic document types

Benefits of Research
This research provides:
1. Improved efficiency in academic document preparation
2. Reduced errors in template processing
3. Enhanced accessibility for academic writers
4. Time savings for institutions and researchers

\chapter{TINJAUAN PUSTAKA}
Previous research in LaTeX document automation has shown various approaches to template processing. Smith et al. (2020) demonstrated the potential of machine learning in document formatting. Jones (2021) explored natural language processing for content generation in academic writing. This study builds upon these foundations while introducing novel approaches to template optimization.

% Original content (commented out):
% \documentclass[9pt, aspectratio=169]{beamer}

\usecolortheme{Blue}
\setbeamersize{text margin left=1.5cm, text margin right=1.5cm}
\useoutertheme{smoothbars}
\useinnertheme{rectangles}

\usepackage[utf8]{inputenc}
\usepackage[brazilian]{babel}
\usepackage{booktabs}
\usepackage{caption}
\usepackage{ragged2e}
\usepackage{subcaption}
\usepackage{tikz}
\usepackage{graphicx}
\usepackage{amsmath}
\usepackage{amsfonts}
\usepackage{amssymb}
\usepackage{epstopdf}
\usepackage{hyperref}
\usepackage{csquotes}
\usepackage[style=numeric]{biblatex}
\usepackage{listings}
\usepackage{FiraSans}
\renewcommand*\familydefault{\sfdefault}


\addbibresource{references.bib}

\setbeamertemplate{frametitle}{
	\vspace{0.5cm}
	\insertframetitle
}


\addbibresource{references.bib}

\usepackage{datetime}
\let\dateBrazilian\relax
\newdateformat{dateBrazilian}{\THEDAY~de \monthname[\THEMONTH]~de \THEYEAR}
\newcommand{\link}[2]{{\color{blue}\href{#2}{#1}}}

\title{Apresentação de Teses e Disestaç\~{o}es}
\subtitle{Uma Visão Geral}
\author{Seu Nome}
\date{\today}


\setbeamertemplate{headline}{} % Remove the blocks from the top panel

%------------------------------------------------------------------------------------------------------

\begin{document}
	\begin{frame}[plain]
		\begin{tikzpicture}[overlay,remember picture]
			\node[anchor=north, inner sep=0.5cm] (logo) at (current page.north)
			{\includegraphics[height=1cm,keepaspectratio]{assets/logo_UCB-PPGE.jpg}};
		\end{tikzpicture}
		\vspace{1cm} % Adjust the vertical space between the logo and the title
		\begin{center}
			\textbf{\Large{\inserttitle}}
			
			\vspace{0.8cm}
			\large{\insertsubtitle}
			
			\vspace{1cm}
			\normalsize{\insertauthor}
			
			\vspace{0.5cm}
			Orientador: Prof. Dr. Nome do Orientador
			
			\vspace{1cm}
			\small{\insertdate}
			
			\vspace{0.5cm}
			\small{Brasília - DF}
		\end{center}
	\end{frame}
	
	%------------------------------------------------------------------------------------------------------
	
	\begin{frame}{Agenda}
		\tableofcontents
	\end{frame}
	
	\section{Introdução}
	% Add your content for the introduction section here
	
	\section{Conceitos Básicos}
	% Add your content for the basic concepts section here
	
	\section{Modelos Econométricos}
	% Add your content for the econometric models section here
	
	\section{Desenvolvimento}
	% Add your content for the econometric methods section here
	
	\section{Resultados}
	% Add your content for the results section here
	
	\section{Conclusão}
	% Add your content for the conclusion section here
	
	%-------------%-------------------------%
	
	\begin{frame}{Introduç\~{a}o}
		%\vspace{1cm}
		% Your frame content here
		Aqui está um exemplo de citação do artigo de Smith e Jones \cite{Author1Title}:
		\begin{quote}
			``Este é um exemplo de citação do artigo de Smith e Jones.''
		\end{quote}
	\end{frame}
	
	\begin{frame}{Modelos Econométricos}
		%\vspace{1cm}
		% Your frame content here
		Aqui está um exemplo de citação da dissertação de mestrado de Eduardo Mangeli \cite{Mangeli2016}:
		\begin{quote}
			``Este é um exemplo de citação da dissertação de mestrado de Eduardo Mangeli.''
		\end{quote}
	\end{frame}
	
	\begin{frame}{Desenvolvimento}
		%\vspace{1cm}
		% Your frame content here
		Aqui está um exemplo de citação do livro de Bernard De Koven \cite{DeKoven2013}:
		\begin{quote}
			``Este é um exemplo de citação do livro de Bernard De Koven.''
		\end{quote}
	\end{frame}
	
	
	\begin{frame}{Resultados}
		%\vspace{1cm}
		% Your frame content here
	\end{frame}
	
	\begin{frame}{Conclus\~{a}o}
		%\vspace{1cm}
		% Your frame content here
	\end{frame}
	
	
	\begin{frame}{Referências}
		\printbibliography[heading=none]
	\end{frame}
	
	\begin{frame}{Perguntas?}
		\begin{tikzpicture}[overlay, remember picture]
			\begin{scope}[opacity=0.7] % Set the opacity for the entire scope
				\node[inner sep=0pt] at (current page.center) {
					\includegraphics[width=\paperwidth, height=\paperheight]{assets/bsb_upscaling(2023).png} % Replace 'background_image.jpg' with your image file
				};
			\end{scope}
		\end{tikzpicture}
		
		\vspace{-3cm} % Adjust vertical space
		\centering
		\Large{Obrigado pela sua atenção!}
		
		\vspace{6pt} % Adjust vertical space
		
		\normalsize{Sinta-se à vontade para fazer perguntas.}
	\end{frame}
	
	
\end{document}


\end{document}